边缘分布很像是正态的, 只有样本 (木板)9 有可能例外 (参见练习 4.33).

为了进一步评价多元正态性, 我们作卡方图, 如\hyperref[fig:4.9]{图 4.9} 所示. 具有最大平方距离的两个样本很清楚偏离了直线图形. 这两个点以及平方距离其次大的一个两点, 使得图形在上面的末端弯曲. 在例 4.15 中我们将回过来再对这一图形进行讨论.

\begin{mvfigure}
\centering
\begin{tikzpicture}[scale=0.62]
\draw[->] (0,0)--(13.2,0) node[right,font=\small]{$q_{c,2}\Bigl(\dfrac{j-\frac{1}{2}}{30}\Bigr)$};
\draw[->] (0,0)--(0,17) node[above,font=\small]{$d_{(j)}^{2}$};
\foreach \x in {2,4,6,8,10,12} \draw (\x,0.15)--(\x,-0.15) node[below,font=\scriptsize]{\x};
\foreach \y in {2,4,6,8,10,12,14,16} \draw (0.15,\y)--(-0.15,\y) node[left,font=\scriptsize]{\y};
% 参考直线 y = x
\draw[dashed] (0,0)--(11,11);
% 数据点（近似 q-q 图）
\foreach \p in {(0.3,0.3),(0.7,0.7),(1.0,0.9),(1.4,1.5),(1.8,1.7),(2.2,2.2),(2.6,2.5),(3.0,3.0),(3.4,3.3),(3.8,3.9),(4.2,4.1),(4.6,4.5),(5.0,5.0),(5.4,5.3),(5.8,5.9),(6.2,6.1),(6.6,6.5),(7.0,7.0),(7.4,7.6),(7.8,8.0),(8.2,8.6),(8.6,9.2),(9.0,9.8),(9.4,10.4),(9.8,11.2)} \fill \p circle (2.2pt);
% 末端偏离点
\fill (10.2,12.0) circle (2.2pt);
\fill (10.6,12.8) circle (2.2pt);
\draw (10.6,12.8) circle (5pt);
\fill (11.0,16.0) circle (2.2pt);
\end{tikzpicture}
\phantomsection\label{fig:4.9}\mvcap{图 4.9\quad 例 4.14 中数据的卡方图}
\end{mvfigure}

我们已经讨论了一些检验正态性假定的简单方法. 特别地, 我们提倡计算 $d_{j}^{2},j=1,2,\cdots,n$ [参见\hyperref[eq:4-32]{式(4-32)}], 并将结果与 $\chi^{2}$ 分位数相比较. 例如, 如果:

1.\quad 大约有一半的 $d_{j}^{2}$ 小于或等于 $q_{c,2}(0.50)$.

2.\quad 分别与 $q_{c,2}\Bigl(\dfrac{1-\frac{1}{2}}{n}\Bigr),q_{c,2}\Bigl(\dfrac{2-\frac{1}{2}}{n}\Bigr),\cdots,q_{c,2}\Bigl(\dfrac{n-\frac{1}{2}}{n}\Bigr)$ 相对应的顺序.

平方距离 $d_{(1)}^{2}\le d_{(2)}^{2}\le\cdots\le d_{(n)}^{2}$ 的图形近似是一条通过原点且斜率为 1 的直线. 那么, 就表明是 $p$ 元正态的 (关于确定正态性方法的更完整的说明参见文献 [6]).

我们在本节结束时, 要注意一切拟合优度的度量方法具有同样严重的缺点. 当样本容量很小, 只有最异常的情况会确定为缺乏拟合. 另一方面, 样本很大时, 总量产生统计上显著的缺乏拟合. 然而, 对规定的分布的偏离可能很小, 这对推断结论在技术上并不重要.

\section{搜寻离群值及“清洁”数据}
\phantomsection\label{sec:4.7}

大多数数据集含有一个或几个异常的观测结果, 它们似乎不属于这个由其他观测值产生的变异性的模型. 在一个单一特性的数据中, 异常观测值相对于其他值来说, 或者是很大或者是很小. 在多元数据下情况可能会更为复杂. 在着手鉴别离群值的问题之前, 必须强调不是所有的离群值都是错误的数据. 它们有可能是一聚类的一部分, 并可以使我们更好地理解所研究的现象.

只要可能, 最好是直观地搜寻出离群值. 当观测值的个数 $n$ 很大时, 点图是不可行的. 当特性 $p$ 的数目很大时, 散布图的数目 $p(p-1)/2$ 很大, 会阻止对它们进行总的观察. 即使如此, 我们建议只要有可能, 应当首先直观地观察数据.

我们应当观察什么? 对于一个单一的随机变量, 问题是一维的, 我们寻找远离其他点的观测值. 例如, 点图
